분할 정복(divide and conquer)은 큰 문제를 작은 문제로 분할한 후 다시 합쳐 문제를 해결하는 방식이다. 분할 정복은 작은 부분에서 구한 결과를 바탕으로 다시 합치는 과정이 간단할 때 유용하다. 이러한 분할 정복은 최적해 문제에 자주 사용되는 특징이 있다. 이번 장에서는 앞에서 배웠던 병합 정렬과, 다른 예시를 통해 분할 정복의 활용성에 대해 알아본다.

\section{다시 보는 병합 정렬}
병합 정렬 과정(\Cref{sec: mergesort})을 다시 보자. 병합 정렬 과정의 핵심을 요약하면 전체 구간을 절반으로 분할한 후 분할된 각 구간이 정렬되어 있음을 이용하여 전체 구간을 $O(N)$의 시간에 정렬하였다. 이러한 방식을 분할 정복 알고리즘이라고 한다. 이러한 알고리즘의 특징 때문에 대부분의 분할 정복 알고리즘은 재귀 함수를 통해 구현된다. \Cref{code: mergesort}를 보면 \texttt{mergesort(st, ed)}가 자기 자신을 다시 호출하는 재귀함수임을 알 수 있다. 

분할 정복은 병합 정렬과 같은 경우 외에도 최적해 문제에서 많이 사용된다. 만약 어떤 구간을 절반으로 분할하여 양쪽에서의 최적값을 찾았다면, 병합 과정에서는 중심을 무조건 포함하는 경우에 대해서만 최적값을 찾으면 되고, 이는 빠른 시간복잡도에 처리할 수 있다. 나중에 배우겠지만, 이러한 분할은 배열 뿐만 아니라 트리에서도 할 수 있어서() 분할 정복 자체는 굉장히 많은 문제와 같이 나올 수 있다고 볼 수 있다.

다음 부분에서 예시 문제인 가장 가까운 두 점을 찾는 문제를 통해 분할 정복으로 최적해 문제를 푸는 방법을 알아보자.

\section{가장 가까운 두 점 찾기}
다음 문제를 보자. \url{https://www.acmicpc.net/problem/2261}과 같다.
\begin{example} \label{ex: nearest-two-points}
    좌표평면상에 주어진 $N$개의 점 중 가장 가까운 두 점 사이의 거리의 제곱을 출력하는 프로그램을 작성하라.
\end{example}

모든 두 점 쌍에 대해 거리를 구하면 거리를 구하는 과정은 상수 시간에 할 수 있으므로 시간복잡도가 $O(N^2)$이 될 것이다. 이제 분할 정복을 이용하여 \textit{like-linear}한 시간에 문제를 해결해보자. 먼저 점을 $x$좌표 순으로 점을 정렬한 후 $x$좌표를 기준으로 좌우로 분할한다. 그리고 왼쪽과 오른쪽에 대해 각각 최단거리를 구한 후 그 중간 부분에 있는 점들의 최단거리를 구할 것이다. 구체적으로 다음과 같이 진행한다.

\begin{itemize}
    \item 점의 개수가 하나이면 리턴한다.
    \item 점의 개수가 두 개 이상이면
    \begin{itemize}
        \item $x$좌표 기준으로 점을 절반으로 나누어 왼쪽 점들 사이에서의 최단거리와 오른쪽 점들 사이에서의 최단거리를 구한다. (함수 호출)
        \item 현재까지의 최단거리 $d$에 대해 경계선 $x=x_c$에서 $d$ 이내의 영역, 즉 $[x_c - d, x_c + d]$에 속하는 점들을 찾고 $y$좌표 기준으로 정렬한다.
        \item 각 점에 대해 $y$좌표 차이가 $d$ 이상이 되기 전까지 비교한다. 이후에는 무조건 거리 차이가 $d$보다 크기 때문에 비교할 필요가 없다.
    \end{itemize}
\end{itemize}

아래는 구현 코드이다.

\begin{code} \label{code: nearest-two-points}
\Cref{ex: nearest-two-points}의 정답 코드
\begin{lstlisting}
#include <bits/stdc++.h>
using namespace std;
typedef long long ll;
#define x first
#define y second
int n;
pair<ll, ll> point[1010101], ylist[1010101], temp[1010101];

bool compare(pair<ll, ll> p1, pair<ll, ll> p2) {
    return p1.y < p2.y;
}
 
ll dis(pair<ll, ll> P, pair<ll, ll> Q) {
    return (P.x - Q.x) * (P.x - Q.x) + (P.y - Q.y) * (P.y - Q.y);
}
 
ll solve(int st, int ed) {
    if (st == ed) {
        ylist[st] = point[st];
        return 9e18;
    }
 
    int mid = (st + ed) / 2;
    ll stdx = point[mid].x, d = 9e18;
    d = min(d, solve(st, mid));
    d = min(d, solve(mid + 1, ed));
 
    sort(ylist + st, ylist + ed + 1, compare);
 
    int r = -1;
    for (int i = st; i <= ed; i++) {
        if ((stdx - ylist[i].x) * (stdx - ylist[i].x) < d) temp[++r] = ylist[i];
    }

    for (int i = 0; i <= r; i++) {
        for (int j = i + 1; j <= r; j++) {
            if ((temp[j].y - temp[i].y) * (temp[j].y - temp[i].y) >= d) break;
            d = min(d, dis(temp[i], temp[j]));
        }
    }
    return d;
}
 
int main() {
    ios_base::sync_with_stdio(false);
    cin.tie(0);
 
    cin >> n;
    for (int i = 1; i <= n; i++) cin >> point[i].x >> point[i].y;
 
    sort(point + 1, point + n + 1);
    cout << solve(1, n) << "\n";
 
    return 0;
}
\end{lstlisting}
\end{code}

\section{마스터 정리}
이제 위 알고리즘의 시간 복잡도를 알아내야 한다. 이를 위해 크기 위 알고리즘의 시간복잡도를 $T(N)$이라고 놓자. 그러면 병합 과정의 시간 복잡도는 정렬이 지배하므로(\Cref{pb: dnc-merge-time}) 다음 식이 성립한다.

\begin{equation*}
    T(N) = \begin{cases} 
        O(1) & \textrm{if } N=1 \\ 
        2T(N/2) + O(N \log N) & \textrm{otherwise} 
     \end{cases}
\end{equation*}

이러한 형태의 식에서 $T(N)$을 구하는 방법은 이미 알려져 있다. 바로 Master Theorem이다.

% \par\noindent\rule{\textwidth}{0.5pt}
% \par

\begin{thm}[Master Theorem]
    다음과 같은 관계식
    \begin{equation*}
        T(n) = \begin{cases} 
            \Theta(1) & \textrm{if } n=1 \\ 
            aT(n/b) + f(n) & \textrm{otherwise} 
         \end{cases}
    \end{equation*}
    에서 $a \geq 1$, $b > 1$이고 $f(n)$이 점근적으로 양의 함숫값을 가지는 함수일 때 어떤 양수 $\varepsilon$ 대해
    \begin{itemize}
        \item $f(n) \in O(n^{\log_b a - \varepsilon}) \Longrightarrow T(n) \in \Theta(n^{\log_b a})$
        \item $f(n) \in \Theta(n^{\log_b a}) \Longrightarrow T(n) \in \Theta(n^{\log_b a} \log n)$
        \item $f(n) \in \Omega(n^{\log_b a + \varepsilon})$, $c<1$과 충분히 큰 $n$에 대해 $af(n/b) \leq cf(n)$ $\Longrightarrow T(n) \in \Theta(f(n))$
    \end{itemize}
    가 성립한다.
\end{thm}

\begin{remark}
    마스터 정리의 두 번째 경우를 확장하면 다음을 얻는다:
    \begin{equation*}
        f(n) \in \Theta(n^{\log_b a} \log^ n) \Longrightarrow T(n) \in \Theta(n^{log_b a} \log^{k+1} n)
    \end{equation*}
\end{remark}

쉽게 설명하면 $O(n^{\log_b a}$보다 $f(n)$의 시간 복잡도가 더 작으면 $T(n)$의 시간복잡도가 $O(n^{\log_b a})$가 되고, 같으면 $O(n^{\log_b a} \log n)$, 크면 $O(f(n))$이 된다. 물론 정확히는 다르지만 재귀적 함수의 시간복잡도를 예측하는데 있어 이 정도의 이해로도 충분할 것이다.

이제 위 알고리즘의 시간복잡도를 마스터 정리로 구해보자. $a=2$, $b=2$, $f(N)=O(N \log N)$이므로 $T(N) = O(N \log^2 N)$이다.

\section{연습문제 A}
\begin{problem} \label{pb: dnc-merge-time}
    \Cref{code: nearest-two-points}의 병합 과정에서 왜 정렬이 시간 복잡도를 지배하려면 병합의 세 번째 과정에서 각 점마다 확인해야 하는 점의 개수가 제한되어야 한다. 실제로 한 점마다 확인해야 하는 점의 개수는 $O(1)$이다. 이를 설명하라.
\end{problem}
\begin{problem}
    사실 \Cref{ex: nearest-two-points}는 $O(N \log N)$에 해결할 수 있다. 해결하는 코드를 작성하고, 실제 백준 문제에 제출해 시간을 비교해보라. (\textit{hint}. 병합 정렬의 아이디어를 활용하라.)
\end{problem}

\section{연습문제 B}

\subsection{기초문제}
\begin{problem} \url{https://www.acmicpc.net/problem/1074} \end{problem}
\begin{problem} \url{https://www.acmicpc.net/problem/1780} \end{problem}
\begin{problem} \url{https://www.acmicpc.net/problem/2630} \end{problem}
\begin{problem} \url{https://www.acmicpc.net/problem/6182} \end{problem}
\begin{problem} \url{https://www.acmicpc.net/problem/6785} \end{problem}
\begin{problem} \url{https://www.acmicpc.net/problem/14600} \end{problem}
\begin{problem} \url{https://www.acmicpc.net/problem/24110} \end{problem}
\begin{problem} \url{https://www.acmicpc.net/problem/24460} \end{problem}

\subsection{응용문제}
\begin{problem} \url{https://www.acmicpc.net/problem/1030} \end{problem}
\begin{problem} \url{https://www.acmicpc.net/problem/1517} \end{problem}
\begin{problem} \url{https://www.acmicpc.net/problem/1588} \end{problem}
\begin{problem} \url{https://www.acmicpc.net/problem/1891} \end{problem}
\begin{problem} \url{https://www.acmicpc.net/problem/2261} \end{problem}
\begin{problem} \url{https://www.acmicpc.net/problem/2339} \end{problem}
\begin{problem} \url{https://www.acmicpc.net/problem/5904} \end{problem}
\begin{problem} \url{https://www.acmicpc.net/problem/8204} \end{problem}
\begin{problem} \url{https://www.acmicpc.net/problem/11219} \end{problem}
\begin{problem} \url{https://www.acmicpc.net/problem/11600} \end{problem}
\begin{problem} \url{https://www.acmicpc.net/problem/14956} \end{problem}
\begin{problem} \url{https://www.acmicpc.net/problem/15278} \end{problem}
\begin{problem} \url{https://www.acmicpc.net/problem/16885} \end{problem}
\begin{problem} \url{https://www.acmicpc.net/problem/18077} \end{problem}
\begin{problem} \url{https://www.acmicpc.net/problem/19808} \end{problem}
\begin{problem} \url{https://www.acmicpc.net/problem/22346} \end{problem}
\begin{problem} \url{https://www.acmicpc.net/problem/24069} \end{problem}

\subsection{도전문제}
\begin{problem} \url{https://www.acmicpc.net/problem/1187} \end{problem}
\begin{problem} \url{https://www.acmicpc.net/problem/3406} \end{problem}
\begin{problem} \url{https://www.acmicpc.net/problem/5482} \end{problem}
\begin{problem} \url{https://www.acmicpc.net/problem/10842} \end{problem}
\begin{problem} \url{https://www.acmicpc.net/problem/15604} \end{problem}
\begin{problem} \url{https://www.acmicpc.net/problem/15842} \end{problem}
\begin{problem} \url{https://www.acmicpc.net/problem/15863} \end{problem}
\begin{problem} \url{https://www.acmicpc.net/problem/16128} \end{problem}
\begin{problem} \url{https://www.acmicpc.net/problem/16349} \end{problem}
\begin{problem} \url{https://www.acmicpc.net/problem/16904} \end{problem}
\begin{problem} \url{https://www.acmicpc.net/problem/17017} \end{problem}
\begin{problem} \url{https://www.acmicpc.net/problem/17679} \end{problem}
\begin{problem} \url{https://www.acmicpc.net/problem/18314} \end{problem}
\begin{problem} \url{https://www.acmicpc.net/problem/18614} \end{problem}
\begin{problem} \url{https://www.acmicpc.net/problem/18621} \end{problem}
\begin{problem} \url{https://www.acmicpc.net/problem/18732} \end{problem}
\begin{problem} \url{https://www.acmicpc.net/problem/18827} \end{problem}
\begin{problem} \url{https://www.acmicpc.net/problem/19635} \end{problem}
\begin{problem} \url{https://www.acmicpc.net/problem/20464} \end{problem}
\begin{problem} \url{https://www.acmicpc.net/problem/20799} \end{problem}
\begin{problem} \url{https://www.acmicpc.net/problem/21543} \end{problem}
\begin{problem} \url{https://www.acmicpc.net/problem/22922} \end{problem}
\begin{problem} \url{https://www.acmicpc.net/problem/27657} \end{problem}